% GENERATED FILE. Edit canonical lessons, then run npm run generate:books.
% canonical-source-hash: fnv1a64:6bcdb6b0a72908d4
% canonical-lessons: TE-C164-turumu, TE-C164-vadakattu, TE-C164-muncu, TE-C164-cimpu, TE-C164-cuttu

\chapter{To grate, To strain, To dip, To tear, To wrap}
\label{ch:te-to-grate-to-strain-to-dip-to-tear-to-wrap}
\hlchaptermodality{164}

\begin{chapteropening}
\textbf{By the end of this chapter:} \emph{I can say to grate, to strain, to dip, to tear and to wrap.}

Complete the last lesson of chapter 164: I can say to grate, to strain, to dip, to tear and to wrap.
\end{chapteropening}

\section[\texorpdfstring{\te{తురుము} (turumu)}{తురుము (turumu)}]{\te{తురుము} (turumu) — to grate}

\label{lesson:TE-C164-turumu}



\begin{quote}
Before the new one: say the Telugu for to compete, then the Telugu for to kick.
\end{quote}

\subsection*{You'll want to know: \te{తురుము}}
\textbf{\te{తురుము}} — \emph{turumu} — \textquotedblleft{}to grate\textquotedblright{}.

\begin{grammarlens}[title={What you've built}]
One more word for this chapter.
\end{grammarlens}

\subsection*{Your turn}
\begin{itemize}
\raggedright
  \item \emph{Say it:} \emph{turumu}
  \item \emph{Say it:} \emph{turumu}, once more
  \item \emph{Say it:} say \emph{turumu}
  \item \emph{Recall:} say the Telugu for to compete, then the Telugu for to kick, then say \emph{turumu} again
\end{itemize}

\begin{tcolorbox}[breakable,colback=teal!4,colframe=teal!35!black,title={Before you move on}]
What does \te{తురుము} mean? (\textquotedblleft{}To grate\textquotedblright{}.) Say it once more.
\end{tcolorbox}
\section[\texorpdfstring{\te{వడకట్టు} (vaḍakaṭṭu)}{వడకట్టు (vaḍakaṭṭu)}]{\te{వడకట్టు} (vaḍakaṭṭu) — to strain, to filter}

\label{lesson:TE-C164-vadakattu}



\begin{quote}
Before the new one: say the Telugu for to kick, then the Telugu for to grate.
\end{quote}

\subsection*{You'll want to know: \te{వడకట్టు}}
\textbf{\te{వడకట్టు}} — \emph{vaḍakaṭṭu} — \textquotedblleft{}to strain, to filter\textquotedblright{}.

\begin{grammarlens}[title={What you've built}]
One more word for this chapter.
\end{grammarlens}

\subsection*{Your turn}
\begin{itemize}
\raggedright
  \item \emph{Say it:} \emph{vaḍakaṭṭu}
  \item \emph{Say it:} \emph{vaḍakaṭṭu}, once more
  \item \emph{Say it:} say \emph{vaḍakaṭṭu}
  \item \emph{Recall:} say the Telugu for to kick, then the Telugu for to grate, then say \emph{vaḍakaṭṭu} again
\end{itemize}

\begin{tcolorbox}[breakable,colback=teal!4,colframe=teal!35!black,title={Before you move on}]
What does \te{వడకట్టు} mean? (\textquotedblleft{}To strain, to filter\textquotedblright{}.) Say it once more.
\end{tcolorbox}
\section[\texorpdfstring{\te{ముంచు} (muñcu)}{ముంచు (muñcu)}]{\te{ముంచు} (muñcu) — to dip, to soak in}

\label{lesson:TE-C164-muncu}



\begin{quote}
Before the new one: say the Telugu for to grate, then the Telugu for to strain.
\end{quote}

\subsection*{You'll want to know: \te{ముంచు}}
\textbf{\te{ముంచు}} — \emph{muñcu} — \textquotedblleft{}to dip, to soak in\textquotedblright{}.

\begin{grammarlens}[title={What you've built}]
One more word for this chapter.
\end{grammarlens}

\subsection*{Your turn}
\begin{itemize}
\raggedright
  \item \emph{Say it:} \emph{muñcu}
  \item \emph{Say it:} \emph{muñcu}, once more
  \item \emph{Say it:} say \emph{muñcu}
  \item \emph{Recall:} say the Telugu for to grate, then the Telugu for to strain, then say \emph{muñcu} again
\end{itemize}

\begin{tcolorbox}[breakable,colback=teal!4,colframe=teal!35!black,title={Before you move on}]
What does \te{ముంచు} mean? (\textquotedblleft{}To dip, to soak in\textquotedblright{}.) Say it once more.
\end{tcolorbox}
\section[\texorpdfstring{\te{చింపు} (cimpu)}{చింపు (cimpu)}]{\te{చింపు} (cimpu) — to tear}

\label{lesson:TE-C164-cimpu}



\begin{quote}
Before the new one: say the Telugu for to strain, then the Telugu for to dip.
\end{quote}

\subsection*{You'll want to know: \te{చింపు}}
\textbf{\te{చింపు}} — \emph{cimpu} — \textquotedblleft{}to tear\textquotedblright{}.

\begin{grammarlens}[title={What you've built}]
One more word for this chapter.
\end{grammarlens}

\subsection*{Your turn}
\begin{itemize}
\raggedright
  \item \emph{Say it:} \emph{cimpu}
  \item \emph{Say it:} \emph{cimpu}, once more
  \item \emph{Say it:} say \emph{cimpu}
  \item \emph{Recall:} say the Telugu for to strain, then the Telugu for to dip, then say \emph{cimpu} again
\end{itemize}

\begin{tcolorbox}[breakable,colback=teal!4,colframe=teal!35!black,title={Before you move on}]
What does \te{చింపు} mean? (\textquotedblleft{}To tear\textquotedblright{}.) Say it once more.
\end{tcolorbox}
\section[\texorpdfstring{\te{చుట్టు} (cuṭṭu)}{చుట్టు (cuṭṭu)}]{\te{చుట్టు} (cuṭṭu) — to wrap, to roll up}

\label{lesson:TE-C164-cuttu}



\begin{quote}
Before the new one: say the Telugu for to dip, then the Telugu for to tear.
\end{quote}

\subsection*{You'll want to know: \te{చుట్టు}}
\textbf{\te{చుట్టు}} — \emph{cuṭṭu} — \textquotedblleft{}to wrap, to roll up\textquotedblright{}.

\begin{grammarlens}[title={What you've built}]
One more word for this chapter.
\end{grammarlens}

\subsection*{Your turn}
\begin{itemize}
\raggedright
  \item \emph{Say it:} \emph{cuṭṭu}
  \item \emph{Say it:} \emph{cuṭṭu}, once more
  \item \emph{Say it:} say \emph{cuṭṭu}
  \item \emph{Recall:} say the Telugu for to dip, then the Telugu for to tear, then say \emph{cuṭṭu} again
\end{itemize}

\begin{tcolorbox}[breakable,colback=teal!4,colframe=teal!35!black,title={Before you move on}]
What does \te{చుట్టు} mean? (\textquotedblleft{}To wrap, to roll up\textquotedblright{}.) Say it once more.
\end{tcolorbox}
